This sections describes the algorithms used by the Value Function Iteration commands. It discusses some of the choices made in the
specific ordering of the commands, as well as how some of the options differ. Covers both the case 1 and case 2 codes, for cpu and gpu.

The focus is on psuedocodes of the commands
\begin{itemize}
 \item \textit{TauchenMethod\_{}Param}: Generates a markov chain approximation for AR(1) process defined by the inputed parameters.
\end{itemize}

\subsection{On the GPU}



\subsubsection{Low Memory}

\subsection{On the CPU}
Implementation on the CPU is similar, but with two main differences. Firstly, rather than giving the Return Function arguement as a
function, it is possible to instead give the Return Function arguement as a matrix. Secondly, the toolkit runs the for-loop over 
$z\_{}c=1:n\_{}z$ (over the z-grid) in parallel on the CPUs (unless $vfoptions.parallel=0$, in which case this is disabled).

\subsection{Howards Improvement Iteration}




